% !TeX root = ../papers/nc-p3.tex
\section{PBW経路格子の行列式による動く族の正則境界の排除}
\label{sec:pbw-valuative-boundaries}
\begin{theorem}\label{thm:pbw-valuative-boundaries}
$K$ を離散付値体とし、標準四窓が $K$ 上で標準Koszul左右AS正則とする。
四位置の文字にそれぞれ全順序を選び、各位置の四文字をその順位で標識する。
位置別左辞書順で六つの二次関係の先頭語が全ての降順対であり、弱い昇順語がPBW基底であるとする。
任意の長さ $n$ の語の正規形の各係数の付値が $-Cn^2$ 以上となる定数 $C$ があるとする。

$W_{i,j}$ を始点 $i$、長さ $j$ のKoszul交叉とし、各々を厳密降順語上の係数で座標化する。
この座標で最上テンソル $w_i$ の唯一の厳密降順語の係数を1に取る。
$w_i\in V_i\otimes W_{i+1,3}$ の四行四列の係数行列の行列式付値を $d_i$ とする。
任意の四位置の基底変更と有限体拡大の後に標準AS正則な特殊ファイバーを持ち得るためには
\[
                         \sum_i d_i\le0
\]
が必要である。これは条件を満たす\emph{全ての}文字順序に課せる。
\end{theorem}
\begin{proof}
各 $W_{i,j}$ の任意の非零元の先頭語は、各隣接対で二次関係の先頭語でなければならない。
実際、関係を置いた位置の外の文字列を固定した各係数列について既約な先頭語を取れる。
従って厳密降順語への射影は単射で、次元 $\binom4j$ が等しいので全単射である。
降順語の単位座標に対応するモニックな基底を取る。全長 $n>0$ のKoszul完全列の各項には
この基底と残りの昇順経路基底との積基底を入れる。この基底での完全列の行列式トーションは $\pm1$ である。
実際、全語の辞書順で各項を濾過すると、次数付き微分で残るのは降順頭部の最後の一文字を
後ろの昇順部へ移しても昇順のままのときの係数1だけである。既約化は全て語を下げ、
モニックな交叉基底の他項も語を下げる。固定全語に寄与する項は、その語が厳密降順の後に
弱い昇順となるときだけで、谷の一文字を左右どちらに含めるかの隣接二項が単位写像で結ばれる。
全語順にこのピボットを除けば主張を得る。

標準正則な格子が存在すると仮定する。その各項の整基底で同じ完全列を取ると、
剰余でも完全なのでトーションは単元である。分数体で二つの列を比較し、
参照PBW座標から整座標への $W_{i,j}$ の行列式付値の四始点の和を
$j=1,2,3,4$ に応じて $V,R,U,M$ と書く。長さ $n$ の経路での和を $b(n)$ とする。
両列のトーション付値が等しいので、定理\ref{thm:regular-orbit-valuations}と同じ
テンソル積の行列式の交代和から
\[
       \sum_{n\ge0} b(n)x^n =
       \frac{P(x)}{(1-x)^8},\qquad P(x)=Vx-Rx^2+Ux^3-Mx^4
\]
を得る。最上四基底は各始点で独立にモニックにしたので、その比較付値も合計 $M$ を使う。
正則格子の $1+3$ 対合の行列式は単元であり、
$(g_i\otimes g_{i+1,3})w_i$ と格子の最上元との倍率を行列式で比較すると
\[
                    V+U+\delta=4M,\qquad\delta=\sum_i d_i .
\]

一次格子の基本因子は両向きに一定範囲にあるから、参照文字の\emph{全語}が生成する
長さ $n$ の格子と正則側の経路格子との差は基本因子ごとに $O(n)$ である。
昇順語格子は全語格子に含まれ、逆には付値の仮定から一様化元の $O(n^2)$ 乗を
全語格子に掛ければ昇順語格子に入る。経路の階数が $\binom{n+3}3$ なので、
ある定数 $C'$ で $-C'n^4\le b(n)\le C'n^5$ を得る。
従って $P(1)=P'(1)=0$ でなければならない。この二式と対合の式を解くと
$R=3M-\delta$, $V=M-\delta/2$, $P''(1)/2=-\delta/2$ となる。
ゆえに $b(n)=-\tfrac12\delta\binom{n+5}{5}+O(n^4)$ であり、下界から $\delta\le0$ を得る。
有限分岐で付値は全て同じ正整数倍となる。
\end{proof}

\begin{corollary}\label{cor:quantum-no-polar-limits}
量子窓の動く族が、任意の基底変更の後に標準AS正則な特殊点を持つなら、
全ての量子三角積の付値は零である。従ってその特殊点は半安定であり、
その半安定商での像は三角積の剰余で定まる量子窓の像に等しい。
\end{corollary}
\begin{proof}
有限拡大後、定理\ref{thm:quantum-component-intersection}の対角twistで各行の交換係数の積を
1に取る。このとき通常の四位置反復がFrobenius標準である。任意の四文字順序を選び、
$j<k$ で $x_kx_j=c_{jk}x_jx_k$ とする。書換えは最大 $n(n-1)/2$ 回で終わる。
降順の最上語から文字 $j$ を先頭に持つ残り降順の語への係数は、符号を除き
$\prod_{k>j}c_{jk}$ である。従って $\delta=4\sum_{j<k}\operatorname{ord}c_{jk}$。
順序反転でこの和は符号を変えるため上の必要条件から零である。さらに任意の隣接二文字だけを
交換した順序と比較すれば各 $\operatorname{ord}c_{jk}=0$ である。
全係数が単元になったので、この量子表示自体が正則な格子を与える。
二つの正則格子の最上テンソルを特殊線形比較に直すと、定理\ref{thm:regular-orbit-valuations}で最上倍率は単元となる。
従って全斉次不変式の特殊値は共通の射影倍率で一致する。
量子窓は多重安定である（24置換語の一様分布で一文字周辺は一様、各Gramは対角）から結論を得る。
\end{proof}

\begin{theorem}\label{thm:all-mixed-flip-component-regular}
片側反転の射影59次元成分の\emph{全ての}標準AS正則点は
$a\in k^\times$ の同型類 $a\sim -a\sim a^{-1}$ で尽きる。
従って正則な同型類の集合は一つのアフィン直線 $a^2+a^{-2}$ であり、
商曲線の無限遠および不安定境界には正則点がない。
\end{theorem}
\begin{proof}
成分の任意の正則点を特殊点に取り、稠密な片側反転族からの曲線と基底変更を、
正規化と有限拡大後の離散付値環上で比較する。$a$ の付値が非零なら逆数同値で正としてよい。
元の文字を $0,1,2,3$ とし、前二文字の標識ビットを位置2,3で反転し、
後二文字の標識ビットを奇位置で反転する。得た標識の数値順を使う。
原稿の再構成から六関係は各辺の六降順対を全て異なる昇順対へ書き換える。
より詳しくは、標識0、標識1、残る二文字の三群の順序を先に整える。
異なる群の入替えは群の組を逆順にするだけで、二文字群内のビットだけが変わり得る。
最後にその群内を整えればよく、全体で $2n^2$ 回以内である。全次数の標準次元から昇順語は基底である。
係数は $a$ の有界な整数冪と四乗根の積なので定理の付値条件も満たす。
\path{computations/verify_pbw_valuative_boundary.py} は、各巡回の
最上語をモニックにし、第一文字と残りの厳密降順三文字の係数行列を直接取る。
四行列式の $a$ の指数は $(2,1,0,-1)$、従って $\delta=2\operatorname{ord}a>0$ となり矛盾する。

残る場合は $a$ が単元なので、片側反転の表示のまま正則に延び、その特殊点は前節で示した閉軌道に属する。
任意の正則模型の特殊点は、二正則格子の比較により半安定で同じ閉軌道を持つ。
この閉軌道の射影次元は全ての $a\ne0$ で58である。実際、上の検算が選ぶ
アフィン軌道微分の59次小行列式は $-64a^{13}(a^2+1)$ であり、
$a=\pm\iota$ でも法109で別の小行列式が非零となる。逆向きの上界は二次元トーラスから従う。
異なる半安定軌道がこれを閉包に含めば、その次元は59以上で、この既約成分で稠密となる。
これは成分上の不変有理関数 $a^2+a^{-2}$ が非定数であることに反する。
従って正則な特殊点は全て表示済みの元の閉軌道に属する。
\end{proof}
